%% -----------------------------------------------------------
\section{Introduction}
%% -----------------------------------------------------------

Mental health disorders represent a leading cause of disability worldwide, yet the populations who bear the greatest burden of psychological distress are systematically excluded from the care systems meant to serve them~\cite{casu_2024}. Sexual and gender minorities experience depressive symptoms at 1.5 times the general population rate, with 59\% of transgender individuals experiencing clinically significant depression and 40\% of \gls{lgbtq} youth seriously considering suicide~\cite{liu_2023}. Blind individuals face elevated rates of depression and anxiety while encountering systematic exclusion from mental health services~\cite{khan_seo_2025}. Neurodivergent individuals, including autistic people and those with \gls{adhd}, experience mental health challenges at substantially elevated rates while finding traditional therapeutic modalities poorly adapted to their cognitive and communicative profiles~\cite{papadopoulos_2025}. Racially and ethnically minoritized communities face documented bias in clinical AI systems that threatens to replicate and amplify existing healthcare disparities~\cite{timmons_2023,haider_2024}. Culturally and linguistically diverse communities struggle with systems that fail to recognize cultural expressions of distress as clinically meaningful~\cite{soubutts_2024}.

AI-powered mental health technologies have emerged with considerable promise into this landscape. They offer on-demand interaction, anonymity, low cost, and no appointment requirements, properties that are particularly valuable for populations for whom stigma and cost constitute primary barriers~\cite{li_2023,drydakis_2022}. Meta-analytic evidence confirms that AI-based conversational agents produce significant reductions in depression (Hedges' $g = 0.64$) and psychological distress (Hedges' $g = 0.70$)~\cite{li_2023,abd_alrazaq_2020}. The field is expanding rapidly: a bibliometric analysis found a 46.19\% annual growth rate in chatbot mental health publications between 2015 and 2024~\cite{zhang_bibliometric_2025}.

Yet rapid expansion has not been accompanied by adequate equity infrastructure. General-purpose chatbots repurposed for mental health have been shown to harm users with depression~\cite{he_2022}. Cultural bias causes AI systems trained on Western data to fail consistently for non-Western populations~\cite{aleem_2024}. Accessibility barriers systematically exclude disabled users from platforms that ostensibly expand access to care~\cite{khan_seo_2025}. Crisis recognition failures place vulnerable users at risk during acute emergencies~\cite{de_freitas_2024}. Demographic bias in \gls{llm}s is now empirically documented at scale across race, gender, and ethnicity~\cite{omar_2025,pfohl_2024}. Socioeconomic barriers compound each of these risks, ensuring that the populations who turn to AI tools out of financial necessity are also those least likely to receive adequate care~\cite{robinson_2024,timmons_2023}. The American Psychological Association has issued explicit guidance warning against generic chatbot use for mental health support~\cite{apa_2025}, and the \gls{who} has called for robust ethics and governance frameworks for large multimodal models in health contexts~\cite{who_2025}.

This survey synthesizes the evidence at this inflection point, where deployment has outpaced safety and equity lags behind scale. It addresses three central research questions:

\begin{itemize}
    \item \textbf{RQ1:} How effective are AI-powered mental health interventions for marginalized populations?
    \item \textbf{RQ2:} What unique risks do marginalized communities face when using AI mental health technologies?
    \item \textbf{RQ3:} Under what conditions can AI mental health technologies be safely and effectively deployed for marginalized populations?
\end{itemize}

The survey contributes by synthesizing evidence across population groups rather than within single identity categories, analyzing equity implications of deployment at scale, identifying design principles grounded in empirical evidence, and surfacing critical research gaps that the field must address as a matter of urgency.